\section{Fused Kernel Compiler}

\label{sec:fusion}

\subsection{Motivation}

A single transformer layer performs approximately 20--30 kernel launches for attention, feed-forward, and normalization operations. Each kernel launch incurs 5--15$\mu$s of overhead on CUDA. For a 32-layer model, this amounts to 3--15ms of pure launch overhead per token---comparable to the actual computation time for small batch sizes.

\subsection{Fusion Opportunities}

We identify three categories of fusable operations:

\begin{table}[H]
\centering
\small
\begin{tabular}{lll}
\toprule
\textbf{Pattern} & \textbf{Operations} & \textbf{Speedup} \\
\midrule
QKV projection & 3 matmuls $\to$ 1 & 2.1x \\
Attention + softmax & Scale, mask, softmax & 1.8x \\
FFN block & Linear, GELU, linear & 1.6x \\
LayerNorm + bias & Norm, scale, shift & 2.4x \\
RoPE embedding & Sin, cos, rotate & 1.9x \\
\bottomrule
\end{tabular}
\caption{Kernel fusion opportunities in transformer architectures.}
\label{tab:fusion}
\end{table}

\subsection{Fusion Algorithm}

The fused kernel compiler operates on the computation graph:

\begin{algorithm}[H]
\caption{Kernel Fusion Pass}
\label{alg:fusion}
\begin{algorithmic}[1]
\Require Computation graph $G = (V, E)$
\State $\text{groups} \gets \emptyset$
\State \Comment{Phase 1: Identify fusable chains}
\For{each node $v \in V$ in topological order}
    \If{$v$ is element-wise or reduction}
        \State Try to merge $v$ into predecessor's group
    \ElsIf{$v$ matches a known pattern (QKV, FFN, etc.)}
        \State Create pattern-specific fused group
    \Else
        \State $v$ is a fusion barrier (start new group)
    \EndIf
\EndFor
\State \Comment{Phase 2: Generate fused kernels}
\For{each group $g \in \text{groups}$}
    \If{$|g| > 1$}
        \State $\text{kernel} \gets \text{CodeGen}(g)$ \Comment{CUDA/Metal kernel}
        \State Replace $g$ with single fused node
    \EndIf
\EndFor
\State \Return optimized graph
\end{algorithmic}
\end{algorithm}

\subsection{Code Generation}

Fused kernels are generated as CUDA PTX or Metal shader code using template-based code generation:

\begin{enumerate}[leftmargin=1.4em]
    \item \textbf{Element-wise fusion}: Chains of unary/binary operations are compiled into a single kernel with loop fusion.
    \item \textbf{Pattern fusion}: Recognized patterns (FlashAttention, fused LayerNorm) map to hand-optimized kernel templates.
    \item \textbf{Reduction fusion}: Reduction operations (softmax, layer norm) are fused with their producers when data locality permits.
\end{enumerate}

\subsection{FlashAttention Integration}

Candle integrates FlashAttention~\cite{dao2022flashattention,dao2023flashattention2} as a native fused operation:

\begin{equation}
\text{Attention}(Q, K, V) = \text{softmax}\left(\frac{QK^T}{\sqrt{d_k}}\right)V,
\end{equation}

implemented as a single fused kernel that:
\begin{enumerate}[leftmargin=1.4em]
    \item Tiles $Q$, $K$, $V$ across shared memory blocks.
    \item Computes attention scores, softmax, and weighted sum in a single pass.
    \item Uses online softmax~\cite{milakov2018online} to avoid materializing the full attention matrix.
    \item Achieves $O(N)$ memory instead of $O(N^2)$ for sequence length $N$.
\end{enumerate}

\subsection{Fusion Results}

Impact of kernel fusion on Llama-3-8B inference (batch size 1, sequence length 2048):

\begin{table}[H]
\centering
\small
\begin{tabular}{lccc}
\toprule
\textbf{Configuration} & \textbf{Kernels} & \textbf{Time (ms)} & \textbf{TPS} \\
\midrule
No fusion & 847 & 42.3 & 23.7 \\
Element-wise only & 612 & 34.1 & 29.3 \\
+ Pattern fusion & 387 & 26.8 & 37.3 \\
+ FlashAttention & 245 & 22.4 & 44.6 \\
\bottomrule
\end{tabular}
\caption{Impact of kernel fusion on inference performance. TPS = tokens per second.}
\label{tab:fusion_results}
\end{table}

